% =============================================================================
% modules/authorization_server.tex — Craftalism Authorization Server
% =============================================================================

\modulelang{Java 17 · Spring Authorization Server · OAuth 2.1 / OIDC}

\chapter{Craftalism Authorization Server}
\label{ch:authorization-server}

\textcolor{CraftGray}{\textit{OAuth 2.1 / OIDC token issuer — machine-to-machine authentication for internal services}}

\bigskip

% =============================================================================
\section{Overview}
% =============================================================================

The Authorization Server is the authentication anchor of the Craftalism platform. It issues short-lived, RSA-signed JWT access tokens to trusted internal clients using the OAuth 2.0 \code{client\_credentials} grant. Downstream services (primarily the Craftalism API) validate these tokens locally by fetching the server's public keys from its JWKS endpoint — no validation round-trip is required per request. The Authorization Server's only registered client is the Minecraft plugin.

% =============================================================================
\section{Responsibilities}
% =============================================================================

\begin{itemize}
  \item Issues RSA-signed JWT access tokens via \code{POST /oauth2/token} using the \code{client\_credentials} grant.
  \item Publishes RSA public keys at \code{GET /oauth2/jwks} for local token verification by downstream services.
  \item Exposes OIDC discovery metadata at \code{GET /.well-known/openid-configuration}.
  \item Exposes OAuth 2.0 server metadata at \code{GET /.well-known/oauth-authorization-server}.
  \item Supports token introspection (\code{/oauth2/introspect}) and revocation (\code{/oauth2/revoke}).
  \item Seeds a single registered OAuth client (\code{minecraft-server} by default) idempotently on every startup.
  \item Enforces a deny-by-default HTTP policy: all paths not explicitly listed as public return 401/403.
\end{itemize}

% =============================================================================
\section{Architecture and Design}
% =============================================================================

The service is a Spring Boot application with two distinct security filter chains and a startup provisioning routine.

\subsection{Filter Chain Model}

Two filter chains run in priority order:

\textbf{Chain 1 — Protocol chain (Order 1):} Handles all \code{/oauth2/**} and discovery endpoints according to OAuth 2.0/OIDC protocol rules. Managed entirely by Spring Authorization Server.

\textbf{Chain 2 — Fallback chain (Order 2):} Stateless and session-less. Permits only the paths listed below; all other requests are denied with \code{401} or \code{403}.

\rowcolors{2}{CraftLight}{white}
\begin{tabularx}{\linewidth}{@{} L{7.5cm} X @{}}
  \toprule
  {\tableheadstyle Permitted path} & {\tableheadstyle Reason} \\
  \midrule
  \code{GET /actuator/health}                       & Health check for orchestrators. \\
  \code{GET /oauth2/jwks}                           & Public key material for JWT verification. \\
  \code{GET /.well-known/oauth-authorization-server} & OAuth 2.0 discovery. \\
  \code{GET /.well-known/openid-configuration}      & OIDC discovery. \\
  \bottomrule
\end{tabularx}

\subsection{Key Configuration}

RSA keys are loaded from environment variables (\code{RSA\_PRIVATE\_KEY} / \code{RSA\_PUBLIC\_KEY}) at startup by \code{RsaKeyConfig}. If neither variable is set, the service generates an ephemeral key pair.

\important{Ephemeral keys are destroyed on service restart. Any JWT issued with an ephemeral key becomes unverifiable after a restart. Never use ephemeral keys in persistent environments.}

\subsection{Client Provisioning}

\code{ClientRegistrationService} seeds one OAuth client (the Minecraft plugin) on every startup. The operation is idempotent: if the client already exists in the database, no change is made. Client credentials, scopes, and grant types are driven by environment variables.

% =============================================================================
\section{Tech Stack}
% =============================================================================

\rowcolors{2}{CraftLight}{white}
\begin{tabularx}{\linewidth}{@{} L{3.0cm} L{5.2cm} X @{}}
  \toprule
  {\tableheadstyle Category} &
  {\tableheadstyle Technology} &
  {\tableheadstyle Notes} \\
  \midrule
  Language      & Java 17                            & \\
  Framework     & Spring Boot 3.5                    & \\
  Security      & Spring Security, Spring Authorization Server & \\
  Database Access & Spring JDBC                      & JDBC repositories for clients/authorizations \\
  Database (runtime) & PostgreSQL                    & \\
  Database (tests)   & H2 (PostgreSQL compat. mode)  & \\
  Build Tool    & Gradle Wrapper                     & \\
  Packaging     & Docker multi-stage image           & \\
  Testing       & JUnit 5, Spring Test               & \\
  \bottomrule
\end{tabularx}

% =============================================================================
\section{Configuration}
% =============================================================================

\rowcolors{2}{CraftLight}{white}
\begin{tabularx}{\linewidth}{@{} L{4.8cm} L{2.6cm} C{1.6cm} X @{}}
  \toprule
  {\tableheadstyle Variable} &
  {\tableheadstyle Default} &
  {\tableheadstyle Required} &
  {\tableheadstyle Description} \\
  \midrule
  \code{DB\_URL}                   & ---                          & \textbf{Yes} & JDBC connection string. \\
  \code{DB\_USER}                  & ---                          & \textbf{Yes} & Database username. \\
  \code{DB\_PASSWORD}              & ---                          & \textbf{Yes} & Database password. \\
  \code{MINECRAFT\_CLIENT\_SECRET} & ---                          & \textbf{Yes} & Secret for the seeded OAuth client. \\
  \code{AUTH\_ISSUER\_URI}         & \code{http://localhost:9000}  & No           & JWT issuer URI. Must match downstream configuration. \\
  \code{MINECRAFT\_CLIENT\_ID}     & \code{minecraft-server}      & No           & Client ID for the seeded OAuth client. \\
  \code{RSA\_PRIVATE\_KEY}         & ---                          & No           & PEM RSA private key (literal \code{\textbackslash n} separators). Ephemeral fallback if absent. \\
  \code{RSA\_PUBLIC\_KEY}          & ---                          & No           & PEM RSA public key (literal \code{\textbackslash n} separators). Ephemeral fallback if absent. \\
  \bottomrule
\end{tabularx}

% =============================================================================
\section{Known Limitations}
% =============================================================================

\begin{itemize}
  \item No \code{generate-keys.sh} script exists in this repository, despite being referenced in documentation and the deployment repository. RSA key pairs must be generated manually.
  \item Integration tests run against H2 in PostgreSQL compatibility mode, not real PostgreSQL. Schema differences may produce false passes.
  \item Token introspection and revocation endpoints have no documented usage examples.
  \item No CI pipeline is configured.
\end{itemize}

% =============================================================================
\section{Roadmap}
% =============================================================================

\begin{itemize}
  \item Add a \code{generate-keys.sh} utility script for RSA key pair generation and formatting.
  \item Add containerized integration tests that run against real PostgreSQL.
  \item Document introspection and revocation usage with curl examples.
  \item Configure a CI pipeline with build and test stages.
\end{itemize}
